%% ============================================================
\section{Authentication, Security, and Access Control}
\label{sec:security}
%% ============================================================

\subsection{Overview}

Security in \echoray{} is implemented as a series of layered controls that
collectively provide defence-in-depth:

\begin{enumerate}[leftmargin=*, label=(\roman*)]
  \item \textbf{Password hashing} at registration (bcrypt, 10 rounds).
  \item \textbf{JWT-based authentication} for stateless request authorisation.
  \item \textbf{Redis token blacklist} for immediate session revocation.
  \item \textbf{Role-based access control} at the service layer.
  \item \textbf{Input validation} via Zod schemas.
  \item \textbf{Transport security} via Helmet middleware and HTTPS/WSS.
\end{enumerate}

\subsection{Password Hashing}

User passwords are hashed using bcrypt (\texttt{utils/auth.ts}) with a cost
factor of~10:

\begin{lstlisting}[caption={Password hashing and verification utilities.},
  label=lst:bcrypt]
import bcrypt from "bcrypt";

export const hashPassword = async (password: string): Promise<string> => {
  return bcrypt.hash(password, 10);
};

export const comparePassword = async (
  password: string,
  hash: string
): Promise<boolean> => {
  return bcrypt.compare(password, hash);
};
\end{lstlisting}

A cost factor of~10 yields approximately~100\,ms hashing time on modern
hardware~\cite{percival2009scrypt}, providing practical brute-force resistance
while maintaining acceptable registration latency.

\subsection{JWT Authentication}

JSON Web Tokens are generated and verified in \texttt{utils/auth.ts}:

\begin{lstlisting}[caption={JWT generation and verification.}, label=lst:jwt]
import jwt from "jsonwebtoken";

export const generateToken = (userId: number): string => {
  return jwt.sign({ userId }, process.env.JWT_SECRET!, {
    expiresIn: "1d",
  });
};

export const verifyToken = (token: string): { userId: number } => {
  return jwt.verify(token, process.env.JWT_SECRET!) as { userId: number };
};
\end{lstlisting}

Tokens are set as \texttt{HttpOnly}, \texttt{Secure}, \texttt{SameSite=Lax}
cookies (preventing client-side JavaScript access and CSRF) and also returned
in the response body for programmatic API consumers and Socket.io handshake
authentication.

\subsection{Token Blacklisting via Redis}

A stateless JWT cannot be revoked before its expiry. \echoray{} resolves this
by maintaining a Redis-backed blacklist (\cref{alg:blacklist}).

\begin{algorithm}[H]
\caption{Token Blacklist Protocol}
\label{alg:blacklist}
\begin{algorithmic}[1]
\Procedure{Logout}{$token$}
  \State $decoded \gets \textsc{VerifyToken}(token)$
  \State $ttl \gets decoded.exp - \lfloor\text{Date.now}()/1000\rfloor$
  \State $\textsc{Redis.setEx}(\texttt{"blacklist:"}+token,\ ttl,\ \texttt{"1"})$
  \State Clear \texttt{token} cookie
\EndProcedure

\Procedure{AuthMiddleware}{$req, res, next$}
  \State $token \gets \textsc{ExtractToken}(req)$
  \If{$token = \emptyset$}
    \State \Return $401$
  \EndIf
  \State $isBlacklisted \gets \textsc{Redis.get}(\texttt{"blacklist:"}+token)$
  \If{$isBlacklisted \ne \emptyset$}
    \State \Return $401$
  \EndIf
  \State $decoded \gets \textsc{VerifyToken}(token)$
  \State $req.user \gets \textsc{GetUserById}(decoded.userId)$
  \State \Call{next}{}
\EndProcedure
\end{algorithmic}
\end{algorithm}

The Redis \texttt{SETEX} command sets both the blacklist entry and its TTL
atomically, guaranteeing that blacklisted tokens are automatically purged when
their original expiry passes, preventing unbounded memory growth.

\subsection{Role-Based Access Control}

\echoray{} implements a two-role model: \textbf{Leader} and \textbf{Member}.
The \texttt{Project.leaderId} field stores the identifier of the user who
created the project; all other collaborators are members.

\begin{definition}[Project Leader]
  A user $u$ is the \emph{leader} of project $p$ if and only if
  $p.\texttt{leaderId} = u.\texttt{id}$.
\end{definition}

\begin{definition}[Privileged Operation]
  An operation $\omega$ is \emph{privileged} if it requires leader status:
  $\omega \in \{\textsc{AddUsers},\, \textsc{RemoveUsers},\, \textsc{DeleteProject}\}$.
\end{definition}

\begin{proposition}[RBAC Enforcement]
  For every privileged operation $\omega$ invoked by user $u$ on project $p$,
  the service layer verifies $p.\texttt{leaderId} = u.\texttt{id}$ and raises
  an authorisation error otherwise.
\end{proposition}

The enforcement code in \texttt{project.service.ts} is representative:

\begin{lstlisting}[caption={RBAC guard in project service.}, label=lst:rbac]
export const addUsersToProject = async (
  projectId: number,
  userIds: number[],
  leaderId: number
): Promise<Project> => {
  const project = await db.project.findUniqueOrThrow({
    where: { id: projectId },
  });
  if (project.leaderId !== leaderId) {
    throw new Error("Only the project leader can add collaborators.");
  }
  return db.project.update({
    where: { id: projectId },
    data: { users: { connect: userIds.map((id) => ({ id })) } },
  });
};
\end{lstlisting}

This service-layer enforcement is complementary to UI restrictions and
guarantees authorisation even in adversarial scenarios where API requests are
crafted directly (e.g., via \texttt{curl}).

\subsection{Input Validation with Zod}

All external inputs are validated with Zod schemas defined in
\texttt{utils/schema.ts}. Selected schemas are shown below.

\begin{lstlisting}[caption={Representative Zod validation schemas
  (\texttt{utils/schema.ts}).}, label=lst:zod]
import { z } from "zod";

export const registerSchema = z.object({
  email:    z.string().email("Invalid email address"),
  password: z.string().min(6, "Password must be at least 6 characters"),
});

export const createProjectSchema = z.object({
  name: z.string()
    .min(1, "Project name is required")
    .max(50, "Project name must be <= 50 characters"),
});

export const addUsersSchema = z.object({
  projectId: z.number().int().positive(),
  userIds:   z.array(z.number().int().positive()).min(1),
});
\end{lstlisting}

Zod's \texttt{safeParse} is used in controllers to return structured
validation errors (HTTP 400) without exposing internal stack traces.

\subsection{HTTP Security Headers}

Helmet~8.1.0 applies a comprehensive set of HTTP security headers:

\begin{itemize}
  \item \texttt{Content-Security-Policy} --- restricts resource loading to
    trusted origins.
  \item \texttt{X-Content-Type-Options: nosniff} --- prevents MIME-type
    sniffing.
  \item \texttt{X-Frame-Options: DENY} --- prevents clickjacking via
    \texttt{<iframe>}.
  \item \texttt{Strict-Transport-Security} --- enforces HTTPS.
  \item \texttt{X-XSS-Protection} --- enables browser XSS filters.
\end{itemize}

CORS is configured with an explicit allow-list
(\texttt{NEXT\_PUBLIC\_CLIENT\_URL}) rather than a wildcard, preventing
cross-origin API abuse.

\subsection{Security Threat Model Summary}

\Cref{tab:threats} maps the STRIDE threat categories~\cite{kohnfelder1999threats}
to \echoray{}'s mitigations.

\begin{table}[H]
  \centering
  \caption{STRIDE threat model and mitigations in \echoray{}.}
  \label{tab:threats}
  \setlength{\tabcolsep}{4pt}
  \begin{tabular}{lll}
    \toprule
    \textbf{Threat}      & \textbf{Example}              & \textbf{Mitigation} \\
    \midrule
    Spoofing             & Forged JWT                    & JWT signature + \texttt{JWT\_SECRET} \\
    Tampering            & SQL injection via prompt      & Prisma parameterised queries \\
    Repudiation          & Deny sending chat message     & Server-side event logging (future) \\
    Info Disclosure      & Password leak                 & bcrypt hashing \\
    Denial of Service    & Flood WebSocket events        & Rate limiting (future) \\
    Elevation of Privilege & Member performing leader ops & RBAC service guard \\
    \bottomrule
  \end{tabular}
\end{table}
